\documentclass[11pt,a4paper]{article}
\usepackage[utf8]{inputenc}
\usepackage{amsmath,amssymb,amsthm}
\usepackage{listings}
\usepackage[margin=2.5cm]{geometry}
\usepackage{xcolor}

\newtheorem{theorem}{Theorem}
\newtheorem{lemma}[theorem]{Lemma}

\lstset{
  language=C++,
  basicstyle=\ttfamily\small,
  keywordstyle=\color{blue}\bfseries,
  commentstyle=\color{green!50!black},
  stringstyle=\color{red!70!black},
  numbers=left,
  numberstyle=\tiny\color{gray},
  breaklines=true,
  frame=single,
  tabsize=4,
  showstringspaces=false
}

\title{IOI 1989 -- Problem 2: Strings}
\author{Solution}
\date{}

\begin{document}
\maketitle

\section{Problem Statement}

Count the number of binary strings of length $n$ that contain no two
consecutive \texttt{1}s.

\section{Solution}

\subsection{Deriving the Recurrence}

Let $f(n)$ denote the count of valid binary strings of length~$n$. Partition
these strings by their last character:
\begin{itemize}
  \item $a(n)$: valid strings of length $n$ ending in \texttt{0}.
  \item $b(n)$: valid strings of length $n$ ending in \texttt{1}.
\end{itemize}

Then $f(n) = a(n) + b(n)$, with recurrences:
\begin{align}
  a(n) &= a(n-1) + b(n-1), \\
  b(n) &= a(n-1).
\end{align}
The first equation holds because \texttt{0} can follow either ending. The
second holds because \texttt{1} can only follow \texttt{0} (to avoid
consecutive \texttt{1}s).

Substituting:
\[
  f(n) = a(n) + b(n) = f(n-1) + a(n-1) = f(n-1) + f(n-2),
\]
where the last step uses $a(n-1) = a(n-2) + b(n-2) = f(n-2)$.

\begin{theorem}
$f(n) = F_{n+2}$, where $F_k$ is the $k$-th Fibonacci number with
$F_1 = F_2 = 1$.
\end{theorem}
\begin{proof}
By induction. We have $f(1) = 2 = F_3$ and $f(2) = 3 = F_4$. For $n \ge 3$,
$f(n) = f(n-1) + f(n-2) = F_{n+1} + F_n = F_{n+2}$.
\end{proof}

\section{Complexity Analysis}

\begin{itemize}
  \item \textbf{Time:} $O(n)$ using iterative computation.
  \item \textbf{Space:} $O(1)$, since only the two most recent values are
    needed.
\end{itemize}

\section{Implementation}

\begin{lstlisting}
#include <iostream>
using namespace std;

int main() {
    int n;
    cin >> n;

    if (n == 0) {
        cout << 1 << endl;  // the empty string
        return 0;
    }
    if (n == 1) {
        cout << 2 << endl;
        return 0;
    }

    long long prev2 = 2;  // f(1)
    long long prev1 = 3;  // f(2)

    for (int i = 3; i <= n; i++) {
        long long cur = prev1 + prev2;
        prev2 = prev1;
        prev1 = cur;
    }

    cout << prev1 << endl;
    return 0;
}
\end{lstlisting}

\section{Verification}

\begin{center}
\begin{tabular}{|c|c|l|}
\hline
$n$ & $f(n)$ & Valid strings \\
\hline
1 & 2 & \texttt{0}, \texttt{1} \\
2 & 3 & \texttt{00}, \texttt{01}, \texttt{10} \\
3 & 5 & \texttt{000}, \texttt{001}, \texttt{010}, \texttt{100}, \texttt{101} \\
4 & 8 & \texttt{0000}, \texttt{0001}, \texttt{0010}, \texttt{0100},
        \texttt{0101}, \texttt{1000}, \texttt{1001}, \texttt{1010} \\
5 & 13 & (Fibonacci pattern continues) \\
\hline
\end{tabular}
\end{center}

The sequence $2, 3, 5, 8, 13, 21, \ldots$ matches $F_3, F_4, F_5, \ldots$,
confirming the formula.

\section{Note on Overflow}

For $n > 85$, the value of $f(n)$ exceeds the range of a 64-bit integer.
If larger values of $n$ are needed, arbitrary-precision arithmetic (e.g.,
Python integers or a big-integer library) should be used.

\end{document}
